\chapter{分子转动、不可约张量与激光诱导排列}\label{ch:molecular-rotation}

双原子分子的取向只需要两个角，而一般刚性多原子分子的构型需要三个 Euler 角。相应地，转动动力学不再只由一个转动惯量控制。本章先从刚体动能建立一般非对称陀螺 Hamiltonian，再说明对称陀螺为什么具有好量子数 \(K\)、非对称陀螺为什么混合 \(K\to K\pm2\)，随后从旋转群的左右作用解释 body-fixed 角动量的“反常”对易关系。最后把分子极化率写成不可约球张量，由 Wigner--Eckart 结构直接推出激光 alignment 的选择定则与时间尺度。

\section{SO(3) 上的量子转子与 Wigner 基}
接下来考察从刚体动能到量子转子 Hamiltonian。

在分子主惯量轴 \(a,b,c\) 中，惯量张量已经对角化：

\[
\mathbf I
=\operatorname{diag}(I_a,I_b,I_c).
\]

经典角速度在 body-fixed 坐标中的分量记为 \(\boldsymbol\Omega=(\Omega_a,\Omega_b,\Omega_c)\)。转动动能和角动量分别为

\[
T_{\rm rot}
=\frac12\boldsymbol\Omega^T\mathbf I\boldsymbol\Omega,
\qquad
J_i=I_i\Omega_i.
\]

消去 \(\Omega_i\) 后得到

\[
T_{\rm rot}
=\frac{J_a^2}{2I_a}
+\frac{J_b^2}{2I_b}
+\frac{J_c^2}{2I_c}.
\]

正则量子化后，刚性多原子转子的 Hamiltonian 为

\begin{equation}
\boxed{
H_{\rm rot}
=\frac{\hat J_a^2}{2I_a}
+\frac{\hat J_b^2}{2I_b}
+\frac{\hat J_c^2}{2I_c}.
}
\label{eq:asymmetric-top-hamiltonian}
\end{equation}

定义能量量纲的主轴转动常数

\[
A_{\rm rot}=\frac{\hbar^2}{2I_a},
\qquad
B_{\rm rot}=\frac{\hbar^2}{2I_b},
\qquad
C_{\rm rot}=\frac{\hbar^2}{2I_c},
\]

并用无量纲角动量 \(\mathbf j=\mathbf J/\hbar\)，则

\begin{equation}
H_{\rm rot}
=A_{\rm rot}j_a^2+B_{\rm rot}j_b^2+C_{\rm rot}j_c^2.
\label{eq:asymmetric-top-energy-form}
\end{equation}

球形陀螺满足 \(I_a=I_b=I_c\)；prolate symmetric top 满足 \(I_a<I_b=I_c\)，即 \(A_{\rm rot}>B_{\rm rot}=C_{\rm rot}\)；oblate symmetric top 满足 \(I_a=I_b<I_c\)，即 \(A_{\rm rot}=B_{\rm rot}>C_{\rm rot}\)。三者都不相等时才是真正的不对称陀螺。

接下来考察Wigner \(D\) 函数与对称陀螺能级。

刚体取向由三个 Euler 角 \((\alpha,\beta,\gamma)\) 描述。旋转群 \(SO(3)\) 上的一组自然正交基是 Wigner 矩阵元

\[
D^J_{MK}(\alpha,\beta,\gamma),
\]

其中 \(J\) 是总角动量量子数，\(M\) 是角动量在 space-fixed \(Z\) 轴上的投影，\(K\) 是在 body-fixed 对称轴上的投影。规范化转动态可记为

\[
\langle\alpha,\beta,\gamma|JKM\rangle
=\sqrt{\frac{2J+1}{8\pi^2}}\,D^{J*}_{MK}(\alpha,\beta,\gamma),
\]

其中归一化采用 Euler 角测度 \(\dd\alpha\,\sin\beta\,\dd\beta\,\dd\gamma\)，且 \(0\le\alpha,\gamma<2\pi\)、\(0\le\beta\le\pi\)。

并满足

\begin{align*}
\mathbf j^2|JKM\rangle&=J(J+1)|JKM\rangle,\\
j_Z|JKM\rangle&=M|JKM\rangle,\\
j_a|JKM\rangle&=K|JKM\rangle
\end{align*}

（最后一式以 prolate 对称轴 \(a\) 为例）。

对于 prolate symmetric top，\eqnrefs{eq:asymmetric-top-energy-form} 化为

\[
H
=B_{\rm rot}\mathbf j^2
+(A_{\rm rot}-B_{\rm rot})j_a^2,
\]

所以 \(|JKM\rangle\) 已经是能量本征态，

\begin{equation}
\boxed{
E_{JK}^{\rm pro}
=B_{\rm rot}J(J+1)
+(A_{\rm rot}-B_{\rm rot})K^2.
}
\label{eq:prolate-symmetric-top-spectrum}
\end{equation}

类似地，oblate symmetric top 取对称轴 \(c\)，得到

\begin{equation}
E_{JK}^{\rm obl}
=B_{\rm rot}J(J+1)
+(C_{\rm rot}-B_{\rm rot})K^2.
\label{eq:oblate-symmetric-top-spectrum}
\end{equation}

无外场时 Hamiltonian 不依赖 space-fixed 投影 \(M\)，因此保持 \(2J+1\) 重简并；对于 \(K\neq0\)，能量又只依赖 \(K^2\)，于是 \(\pm K\) 成对简并。


\begin{derivation}[Wigner $D$ 基底的正交性来自 Haar 测度与 Schur 引理，即\eqnrefs{eq:wigner-d-haar-orthogonality}]
\emph{设置。}令 $D^J(R)$ 是 $SO(3)$ 的不可约酉表示，维数 $2J+1$，以归一化 Haar 测度 $\dd R$ 满足 $\int_{SO(3)}\dd R=8\pi^2$。Haar 测度的关键性质是左右不变：$\int f(R_0R)\dd R=\int f(R)\dd R$ 对任意 $R_0$。这一不变性是紧群的普遍结构，不依赖 $SO(3)$ 的具体参数化。在 Euler 角参数化 $R=R(\alpha,\beta,\gamma)$ 下，Haar 测度显式为 $\dd R=\dd\alpha\,\sin\beta\,\dd\beta\,\dd\gamma$，积分区域 $0\le\alpha,\gamma<2\pi$、$0\le\beta\le\pi$，总体积 $\int\dd R=(2\pi)(2)(2\pi)=8\pi^2$。$\sin\beta$ 因子来自 $SO(3)$ 作为 $S^3/\mathbb Z_2$ 的自然体积元，保证左右不变性。

\emph{第一步：构造 intertwining 算符。}固定矩阵元指标，考虑算符
\begin{equation*}
 X
 =\int\dd R\,
 D^J(R)\,|K\rangle\langle K'|\,D^{J'}(R)^\dagger.
\end{equation*}
这是从 $J'$ 表示空间到 $J$ 表示空间的线性映射。$X$ 把 $J'$ 表示的基矢 $|K'\rangle$ 通过群平均"搬运"到 $J$ 表示空间。群平均是构造 intertwining 算符的标准技巧：对任意表示 $\rho$，$\int\dd R\,\rho(R)A\rho'(R)^\dagger$ 自动满足 intertwining 条件，Haar 不变性是关键。

\emph{第二步：Haar 不变性给出 intertwining。}对任意 $R_0\in SO(3)$，由 Haar 测度左不变性，
\begin{equation*}
 D^J(R_0)X
 =\int\dd R\,D^J(R_0R)|K\rangle\langle K'|D^{J'}(R)^\dagger
 =\int\dd R'\,D^J(R')|K\rangle\langle K'|D^{J'}(R_0^{-1}R')^\dagger
 =XD^{J'}(R_0),
\end{equation*}
其中第二步作了换元 $R'=R_0R$，第三步用了 $D^{J'}(R_0^{-1}R')^\dagger=D^{J'}(R')^\dagger D^{J'}(R_0)$。因此 $X$ 是 $D^{J'}$ 到 $D^J$ 的 intertwining 算符，满足 $D^J(R_0)X=XD^{J'}(R_0)$ 对所有 $R_0$。注意只用了 Haar 测度的左不变性；右不变性会给出 $X^\dagger$ 的类似 intertwining 条件，两者结合保证 $X$ 的结构同时受左右约束。

\emph{第三步：Schur 引理定 $X$ 的结构。}若 $J\neq J'$，两个不可约表示不等价，Schur 引理要求 $X=0$。若 $J=J'$，$X$ 与不可约表示自身交换，Schur 引理要求 $X$ 正比于单位算符：$X=c\,\delta_{KK'}\,I$。Schur 引理的证明：若 $X$ 与所有 $D^J(R)$ 交换，取 $D^J$ 的任一本征基，$X$ 在该基中是对角的；但不可约表示的所有矩阵元在酉等价意义下遍历整个 $\mathrm{U}(2J+1)$，迫使 $X$ 的所有本征值相等，即 $X=cI$。取矩阵元 $\langle M|X|M'\rangle$，利用 $D^J(R)^\dagger=D^J(R^{-1})$ 和表示的酉性，
\begin{equation}\label{eq:wigner-d-haar-orthogonality}
 \int\dd R\,
 D^{J*}_{MK}(R)\,D^{J'}_{M'K'}(R)
 =\frac{8\pi^2}{2J+1}
 \delta_{JJ'}\delta_{MM'}\delta_{KK'},
\end{equation}
即\eqnrefs{eq:wigner-d-haar-orthogonality}。归一化常数 $8\pi^2/(2J+1)$ 来自 $\int\dd R=8\pi^2$ 除以表示维数 $2J+1$。三个 Kronecker delta 分别来自：$\delta_{JJ'}$（不等价表示间 $X=0$）、$\delta_{MM'}$（$X\propto I$ 在 $M$ 基中对角）、$\delta_{KK'}$（构造 $X$ 时已固定 $K,K'$）。

\emph{Schur 引理的两种形式。}本推导用到 Schur 引理的两个等价表述：(i) 不同不可约表示间的 intertwining 算符为零；(ii) 不可约表示到自身的 intertwining 算符为一维（仅恒等算子的标量倍）。两者都建立在表示不可约性上：若表示可约，$X$ 可在不同不可约分量间非零，正交关系会含额外指标。$SO(3)$ 的不可约表示由 $J$ 唯一标记，故正交关系只含 $J,J',M,M',K,K'$ 六个指标。

\emph{归一化因子的几何含义。}$\sqrt{(2J+1)/(8\pi^2)}$ 来自 $SO(3)$ 的体积 $8\pi^2$ 与表示维数 $2J+1$ 之比。在 Euler 角参数化下，$8\pi^2=(2\pi)^2\cdot2$ 来自 $\alpha,\gamma$ 各贡献 $2\pi$、$\beta$ 贡献 $2$（$\int_0^\pi\sin\beta\,\dd\beta=2$）。这一因子保证 $\langle JKM|JKM\rangle=1$，即转动态在 $SO(3)$ 上的概率密度归一化。

\emph{左右作用的对偶性。}正交关系中 $D^{J*}_{MK}$ 和 $D^{J'}_{M'K'}$ 分别涉及 space-fixed 指标 $M$ 和 body-fixed 指标 $K$。$M$ 的正交来自左作用（主动旋转），$K$ 的正交来自右作用（被动旋转）。两种作用的独立性（$[L_i,R_j]=0$）保证 $M$ 和 $K$ 可同时正交化，这是 $|JKM\rangle$ 三指标正交的群论根源。这一对偶性是紧群表示论的普遍特征：左正则表示与右正则表示在 $L^2(G)$ 上交换，各自分解为相同的不可约表示。

\emph{第四步：完备性。}Peter--Weyl 定理进一步保证这些矩阵元在 $L^2(SO(3))$ 中完备，因此一般刚体波函数都可按 $|JKM\rangle$ 展开；对称陀螺之所以额外保留好量子数 $K$，是因为 Hamiltonian 还与对应的右作用生成元对易。展开系数
\begin{equation*}
 \Psi(\alpha,\beta,\gamma)=\sum_{JKM}c_{JKM}\sqrt{\frac{2J+1}{8\pi^2}}\,D^{J*}_{MK}(\alpha,\beta,\gamma),
\end{equation*}
由正交关系给出 $c_{JKM}=\int\dd R\,\Psi(R)\sqrt{(2J+1)/(8\pi^2)}\,D^J_{MK}(R)$。对称陀螺的 Hamiltonian $H=B_{\rm rot}\mathbf j^2+(A_{\rm rot}-B_{\rm rot})j_a^2$ 与 $j_a^2$ 对易，故 $K$ 是好量子数，展开只在固定 $K$ 内进行；不对称陀螺的 $j_\pm^2$ 项破坏 $K$ 守恒，需要在固定 $J,M$ 的 $(2J+1)$ 维空间内对角化混合矩阵。

\emph{举例：$J=0$ 和 $J=1$ 的正交核对。}$J=0$ 表示是常数 $D^0_{00}=1$，$\int\dd R\,|D^0_{00}|^2=8\pi^2$，与 $8\pi^2/(2\cdot0+1)=8\pi^2$ 一致。$J=1$ 表示是 $3\times3$ 旋转矩阵 $R_{ij}$，$\int\dd R\,R_{ij}R_{kl}=\frac{8\pi^2}{3}\delta_{ik}\delta_{jl}$，9 个矩阵元彼此正交，每个的模平方平均为 $8\pi^2/3$。这直接验证了正交关系。在 Euler 角参数化 $(\alpha,\beta,\gamma)$ 下，$D^1_{00}=P_1(\cos\beta)=\cos\beta$，$\int_0^\pi\cos^2\beta\sin\beta\,\dd\beta\cdot(2\pi)^2=8\pi^2/3$，与一般公式一致。

\emph{与不对称陀螺的关系。}不对称陀螺 Hamiltonian \eqnrefs{eq:asymmetric-top-k-mixing} 在 $|JKM\rangle$ 基中是有限矩阵，正交关系保证该基在固定 $J,M$ 的 $(2J+1)$ 维子空间中完备。不对称性通过 $j_\pm^2$ 项混合 $K\to K\pm2$，但正交基保证了矩阵元 \eqnrefs{eq:asymmetric-top-offdiagonal-element} 的解析表达式。Wigner $D$ 函数的正交性因此是不对称陀螺能级可数值对角化的数学基础。

\emph{更一般结论：非紧群与 Peter--Weyl。}Peter--Weyl 定理对任意紧群成立：$L^2(G)$ 分解为所有不可约表示矩阵元的直和，每个不可约表示 $D^\lambda$ 出现 $\dim(\lambda)$ 次。对 $SO(3)$，不可约表示由 $J$ 标记，$\dim=2J+1$，故 $L^2(SO(3))=\bigoplus_J(2J+1)\mathcal H_J$。对非紧群（如 $SO(3,1)$ 或 Heisenberg 群），不可约表示一般是无穷维的，正交关系由 Plancherel 测度替代 Haar 测度的离散求和。分子转动的 $SO(3)$ 框架因此是紧群调和分析的特例；若考虑分子在非紧对称空间上的运动（如刚性转子在双曲平面上的扩散），需要相应的非紧群表示论。

\emph{举例：$J=2$ 的正交维度检验。}$J=2$ 表示维数为 5，正交关系给出 $\int\dd R\,D^{2*}_{MK}D^2_{M'K'}=\frac{8\pi^2}{5}\delta_{MM'}\delta_{KK'}$，25 个矩阵元彼此正交，每个模平方为 $8\pi^2/5\approx15.8$。这比 $J=1$ 的 $8\pi^2/3\approx26.3$ 更小，反映了高 $J$ 表示的矩阵元在群上更"弥散"——维数越高，每个矩阵元携带的 $L^2$ 范数越小。

\emph{与球谐函数的关系。}当 $K=0$ 时，$D^J_{M0}(\alpha,\beta,\gamma)=\sqrt{4\pi/(2J+1)}\,Y_{JM}^*(\beta,\alpha)$ 与 $\gamma$ 无关，正交关系退化为球谐函数的正交性 $\int Y_{JM}^*Y_{J'M'}\,\dd\Omega=\delta_{JJ'}\delta_{MM'}$。这把 Wigner $D$ 正交关系与 $S^2$ 上的调和分析联系起来：球谐函数是 $SO(3)/SO(2)\cong S^2$ 上的诱导表示。

\emph{物理意义。}Wigner \(D\) 矩阵正交性的物理核心是"不可约表示在群上的正交性"：不同不可约表示的矩阵元在 Haar 测度下正交，同一表示的不同矩阵元也正交，归一化因子 \(8\pi^2/(2J+1)\) 来自 Schur 引理给出的 intertwiner 维数。物理后果包括：(1) 任意群函数可按 \(D^J_{MK}\) 展开为"群 Fourier 级数"——这是刚性转子波函数 \(\psi(\alpha,\beta,\gamma)=\sum_{JMK}c_{JMK}D^J_{MK}\) 的理论基础，系数 \(c_{JMK}\) 是群 Fourier 变换；(2) 正交关系保证不同 \(J\) 的转动能级非简并耦合，刚性转子的能级 \(E_J=BJ(J+1)\) 来自 Casimir 算子 \(J^2\) 的本征值；(3) 分子微波光谱中转动跃迁 \(\Delta J=\pm1\)、\(\Delta K=0\) 的选择定则来自 \(D^J_{MK}\) 在偶极算符（rank-1 张量）下的矩阵元结构。\(K=0\) 退化为球谐函数给出轨道角动量量子力学，\(K\ne0\) 的完整 \(D\) 矩阵描述对称陀螺分子（如 NH\(_3\)）的转动波函数。
\end{derivation}

\section{左右生成元与不对称陀螺}
接下来考察不对称陀螺的 \(K\) 混合。

选择 body-fixed \(c\) 轴作为量子化轴，定义无量纲升降算符

\[
j_\pm=j_a\pm\ii j_b.
\]

利用

\[
j_a=\frac{j_++j_-}{2},
\qquad
j_b=\frac{j_+-j_-}{2\ii},
\]

有

\begin{align*}
A_{\rm rot}j_a^2+B_{\rm rot}j_b^2
&=\frac{A_{\rm rot}+B_{\rm rot}}{4}
(j_+j_-+j_-j_+)\\
&\quad+\frac{A_{\rm rot}-B_{\rm rot}}{4}
(j_+^2+j_-^2).
\end{align*}

再用

\[
\frac12(j_+j_-+j_-j_+)=\mathbf j^2-j_c^2,
\]

可将一般不对称陀螺写成

\begin{equation}
H_{\rm rot}
=\frac{A_{\rm rot}+B_{\rm rot}}{2}\mathbf j^2
+\left(C_{\rm rot}-\frac{A_{\rm rot}+B_{\rm rot}}2\right)j_c^2
+\frac{A_{\rm rot}-B_{\rm rot}}4(j_+^2+j_-^2).
\label{eq:asymmetric-top-k-mixing}
\end{equation}

在 \(|JKM\rangle\) 基中，对角矩阵元为

\begin{equation}
\langle JKM|H_{\rm rot}|JKM\rangle
=\frac{A_{\rm rot}+B_{\rm rot}}2J(J+1)
+\left(C_{\rm rot}-\frac{A_{\rm rot}+B_{\rm rot}}2\right)K^2.
\label{eq:asymmetric-top-diagonal-element}
\end{equation}

而

\[
j_+|JKM\rangle
=\sqrt{J(J+1)-K(K+1)}\,|J,K+1,M\rangle,
\]

连续作用两次得到

\begin{align*}
j_+^2|JKM\rangle
&=C^+_{JK}|J,K+2,M\rangle,\\
C^+_{JK}
&=\sqrt{[J(J+1)-K(K+1)]
[J(J+1)-(K+1)(K+2)]}.
\end{align*}

因此唯一的非零带外矩阵元是

\begin{equation}
\boxed{
\langle J,K+2,M|H_{\rm rot}|JKM\rangle
=\frac{A_{\rm rot}-B_{\rm rot}}4 C^+_{JK},
}
\label{eq:asymmetric-top-offdiagonal-element}
\end{equation}

以及其 Hermitian 共轭 \(K\to K-2\)。所以非对称性不是任意混合所有 \(K\)，而是把固定 \(J,M\) 子空间分裂成若干只通过 \(\Delta K=\pm2\) 相连的块。\(K\) 不再是好量子数，但 \(K\) 的奇偶性等离散对称信息仍然约束矩阵结构。

常用 Ray asymmetry parameter 记为

\begin{equation}
\boxed{
\kappa_{\rm R}
=\frac{2B_{\rm rot}-A_{\rm rot}-C_{\rm rot}}
{A_{\rm rot}-C_{\rm rot}},
\qquad -1\le\kappa_{\rm R}\le1.
}
\label{eq:ray-asymmetry-parameter}
\end{equation}

其中 \(\kappa_{\rm R}=-1\) 是 prolate 极限，\(\kappa_{\rm R}=+1\) 是 oblate 极限。一般 \(\kappa_{\rm R}\) 下没有单个闭式 \(E_{JK}\)，应在每个固定 \(J,M\) 的有限矩阵中对角化 \eqnrefs{eq:asymmetric-top-k-mixing}。

接下来考察body-fixed 角动量对易关系的群论来源。

同一个转子同时允许两种群作用。space-fixed 转动改变分子在实验室中的整体取向，可以看成从左侧作用于旋转矩阵；body-fixed 转动则改变分子坐标架的内部标记，对应右侧作用。设 \(L_i\) 与 \(R_i\) 分别是 \(SO(3)\) 上的左不变和右不变向量场。Lie 群上的一般结果是

\begin{equation}
[L_i,L_j]=\varepsilon_{ijk}L_k,
\qquad
[R_i,R_j]=-\varepsilon_{ijk}R_k,
\qquad
[L_i,R_j]=0.
\label{eq:left-right-invariant-algebra}
\end{equation}

负号来自右不变场的群交换子顺序与 Lie 代数元素乘法顺序相反。量子角动量由 \(-\ii\hbar\) 乘相应微分生成元得到，因此

\begin{equation}
[J_X,J_Y]=+\ii\hbar J_Z,
\qquad
[J_a,J_b]=-\ii\hbar J_c,
\label{eq:body-fixed-anomalous-commutator}
\end{equation}

以及循环置换。

\begin{derivation}[右作用为什么产生负号]
分两步比较左不变与右不变生成元的交换子。核心是群回路的乘积顺序在左右作用下的差异。

\emph{第一步：左不变场的交换子。}取 Lie 代数元素 \(X,Y\)，小参数群元 \(\ee^{\epsilon X}\)、\(\ee^{\eta Y}\)。左不变生成元 \(L_X\) 在群元 \(g\) 处的切向由 \(\ee^{\epsilon X}g\) 定义，即群从左侧作用。群回路
\[
 \ee^{\epsilon X}\ee^{\eta Y}\ee^{-\epsilon X}\ee^{-\eta Y}
\]
从恒元出发，先右乘 \(\ee^{-\eta Y}\)，再右乘 \(\ee^{-\epsilon X}\)，等等。由 BCH 公式逐阶展开：
\[
 \ee^{\epsilon X}\ee^{\eta Y}
 =\exp\!\left(\epsilon X+\eta Y+\frac{\epsilon\eta}{2}[X,Y]+O(\epsilon^2\eta,\epsilon\eta^2)\right),
\]
\[
 \ee^{\epsilon X}\ee^{\eta Y}\ee^{-\epsilon X}
 =\exp\!\left(\eta Y+\epsilon\eta[X,Y]+O(\epsilon^2\eta,\epsilon\eta^2)\right),
\]
\[
 \ee^{\epsilon X}\ee^{\eta Y}\ee^{-\epsilon X}\ee^{-\eta Y}
 =\exp\!\left(\epsilon\eta[X,Y]+O(\epsilon^2\eta,\epsilon\eta^2)\right).
\]
回路的一阶项 $\epsilon X$、$\eta Y$ 逐对消去，二阶项留下 $[X,Y]$。因此 \([L_X,L_Y]=L_{[X,Y]}\)，与 Lie 代数的原始交换子符号一致。

\emph{第二步：右不变场的回路反向。}右不变生成元 \(R_X\) 在群元 \(g\) 处的切向由 \(g\ee^{\epsilon X}\) 定义，即群从右侧作用。同样的微分算符组合 \(R_XR_Y-R_YR_X\) 作用在函数上时，群元以相反次序累积：先左乘 \(\ee^{\epsilon X}\)，再左乘 \(\ee^{\eta Y}\)，故对应回路为
\[
 \ee^{-\eta Y}\ee^{-\epsilon X}\ee^{\eta Y}\ee^{\epsilon X}.
\]
与第一步的回路相比，这是逆元置换 \(\epsilon\to-\epsilon\)、\(\eta\to-\eta\)，因此 BCH 展开给出
\[
 \ee^{-\eta Y}\ee^{-\epsilon X}\ee^{\eta Y}\ee^{\epsilon X}
 =\exp\!\left(-\epsilon\eta[X,Y]+O(\epsilon^2\eta,\epsilon\eta^2)\right).
\]
因此 \([R_X,R_Y]=-R_{[X,Y]}\)。

\emph{第三步：量子化的符号约定。}量子角动量由 $J_i=-\ii\hbar L_i$（space-fixed）或 $J_a=-\ii\hbar R_a$（body-fixed）定义。代入交换子，
\begin{equation*}
 [J_X,J_Y]=(-\ii\hbar)^2[L_X,L_Y]=-\hbar^2 L_{[X,Y]}=+\ii\hbar J_Z,
\end{equation*}
\begin{equation*}
 [J_a,J_b]=(-\ii\hbar)^2[R_a,R_b]=-\hbar^2(-R_{[a,b]})=+\hbar^2 R_{[a,b]}=-\ii\hbar J_c,
\end{equation*}
即\eqnrefs{eq:body-fixed-anomalous-commutator}。负号来自右不变场的负交换子，而非量子化异常。

\emph{物理含义。}所谓"反常"只是在同一个 \(SO(3)\) 上选择右不变而不是左不变生成元的结果：右不变场对应于被动变换（坐标旋转），而左不变场对应于主动变换（态旋转），二者差一个整体符号，并不表示物理角动量违反了旋转代数。在分子物理中，space-fixed 角动量 $J_X,J_Y,J_Z$ 描述实验室系中分子整体取向的变化，body-fixed 角动量 $J_a,J_b,J_c$ 描述分子内坐标架的标记变化。两组生成元彼此对易 $[J_i,J_a]=0$（\eqnrefs{eq:left-right-invariant-algebra} 第三式），故可同时标记量子数 $M$（space-fixed $Z$ 投影）和 $K$（body-fixed 对称轴投影）。

\emph{举例：对称陀螺好量子数 $K$ 的群论根源。}对 prolate symmetric top，Hamiltonian $H=B_{\rm rot}\mathbf j^2+(A_{\rm rot}-B_{\rm rot})j_a^2$ 与 $j_a$ 对易（$j_a$ 是右作用生成元），故 $K$ 是好量子数。不对称陀螺的 $j_\pm^2$ 项与 $j_a$ 不对易但与 $j_a^2$ 对易，故 $K$ 不再好但 $K^2$（即 $K$ 的奇偶性）仍守恒。这一结构完全由左右作用的群论决定，与具体分子参数无关。例如 $J=2$ 的不对称陀螺有 $K=0,\pm1,\pm2$ 五个态，$K$ 奇偶性把 $5\times5$ 矩阵分成 $K\in\{0,\pm2\}$（3 维）和 $K\in\{\pm1\}$（2 维）两个块，各自独立对角化。

\emph{不对称陀螺的 Ray 参数极限。}$\kappa_{\rm R}=-1$（prolate 极限）时 $j_\pm^2$ 项消失，$K$ 恢复为好量子数；$\kappa_{\rm R}=+1$（oblate 极限）同理。$\kappa_{\rm R}=0$（最大不对称）时 $K$ 混合最强，能级对 $\kappa_{\rm R}$ 的依赖在分子转动光谱中表现为不对称分裂图样，是鉴定分子构型的实验依据。水分子（$C_{2v}$，$\kappa_{\rm R}\approx-0.44$）和甲醛（$C_{2v}$，$\kappa_{\rm R}\approx-0.96$）的不对称参数差异直接反映其惯量主轴比的差异。

\emph{更一般结论：非紧 Lie 群的右作用。}对非紧半单 Lie 群（如 Lorentz 群 $SO(3,1)$），左不变场仍满足 $[L_X,L_Y]=L_{[X,Y]}$，但右不变场的负号关系 $[R_X,R_Y]=-R_{[X,Y]}$ 同样成立。负号的根源是群的乘法顺序反转，与紧性无关。在规范理论中，这一符号差异出现在规范势的左右作用约定中，是 Yang--Mills 理论中结构常数符号约定的 Lie 群论基础。
\end{derivation}

\section{不可约极化率张量与 alignment 选择定则}
接下来考察诱导偶极相互作用的 rank-0/rank-2 分解。

对远离真实电子共振、且电场频率足够高到可以对光学周期平均的非极性分子，主要相互作用来自诱导偶极

\[
\mathbf p(t)=\boldsymbol\alpha\mathbf E(t),
\]

其中 \(\boldsymbol\alpha\) 是分子极化率二阶张量。若实电场写成 \(\mathbf E(t)=\mathbf E_0\cos\omega t\)，周期平均相互作用为

\begin{equation}
\boxed{
\overline V_{\rm int}
=-\frac14\mathbf E_0^T\boldsymbol\alpha\mathbf E_0.
}
\label{eq:polarizability-cycle-averaged}
\end{equation}

对线性分子

\[
\boldsymbol\alpha
=\alpha_\perp\mathbf 1
+\Delta\alpha\,\hat{\mathbf n}\hat{\mathbf n},
\qquad
\Delta\alpha=\alpha_\parallel-\alpha_\perp.
\]

若线偏振沿实验室 \(Z\) 轴，\(\theta\) 是分子轴和 \(Z\) 的夹角，则

\begin{equation}
V(\theta)
=-\frac14E_0^2
[\alpha_\perp+\Delta\alpha\cos^2\theta].
\label{eq:linear-molecule-alignment-potential}
\end{equation}

利用

\[
\cos^2\theta
=\frac13+\frac23P_2(\cos\theta),
\]

可把 \eqnrefs{eq:linear-molecule-alignment-potential} 分成

\begin{align*}
V(\theta)
&=-\frac14E_0^2
\left(\alpha_\perp+\frac{\Delta\alpha}{3}\right)
-\frac{E_0^2\Delta\alpha}{6}P_2(\cos\theta).
\end{align*}

第一项是 rank-0 标量，只给所有转动态共同的 AC Stark 位移；真正改变转动态组成的是第二个 rank-2 项。这就是为什么 alignment 的角动量选择定则由二阶球张量而不是一阶偶极张量决定。

一般地，两个 rank-1 场张量的耦合写成

\[
[E^{(1)*}\otimes E^{(1)}]^{(K)}_Q,
\qquad K=0,1,2.
\]

对非磁、无手性且极化率对称的普通分子，反对称 rank-1 部分消失，保留 \(K=0,2\)。两个同阶不可约张量的旋转标量积定义为

\[
T^{(K)}\!\cdot U^{(K)}
=\sum_{Q=-K}^{K}(-1)^Q T^{(K)}_QU^{(K)}_{-Q},
\]

于是相互作用可以统一写成

\begin{equation}
H_{\rm int}
=-\frac12\sum_{KQ}(-1)^Q
[E^{(1)*}\otimes E^{(1)}]^{(K)}_Q
\alpha^{(K)}_{-Q},
\qquad K=0,2,
\label{eq:polarizability-tensor-coupling}
\end{equation}

其中整体数值因子随复振幅定义改变，但张量结构与选择定则不变。

接下来考察从 Gaunt 积分推出 alignment 的选择定则。

对线性刚性转子，场沿 \(Z\) 时只有 \(P_2(\cos\theta)\) 需要计算。利用

\[
P_2(\cos\theta)=\sqrt{\frac{4\pi}{5}}Y_{20}(\theta,\phi)
\]

和三个球谐函数的 Gaunt 积分，可以得到

\begin{equation}
\begin{aligned}
&\langle J'M'|P_2(\cos\theta)|JM\rangle\\
&\quad=(-1)^{M'}
\sqrt{(2J'+1)(2J+1)}
\begin{pmatrix}
J'&2&J\\0&0&0
\end{pmatrix}
\begin{pmatrix}
J'&2&J\\-M'&0&M
\end{pmatrix}.
\end{aligned}
\label{eq:alignment-p2-matrix-element}
\end{equation}

第二个 3-j 符号要求

\[
M'=M,
\]

三角条件要求 \(|J-J'|\le2\le J+J'\)，而第一个 3-j 符号还要求 \(J+J'+2\) 为偶数。因此线偏振 alignment 的主要选择定则是

\begin{equation}
\boxed{
\Delta M=0,
\qquad
\Delta J=0,\pm2.
}
\label{eq:alignment-selection-rule}
\end{equation}

这与电偶极跃迁常见的 \(\Delta J=\pm1\) 根本不同，因为这里驱动系统的是 rank-2 极化率各向异性。

\section{绝热、脉冲 kick 与转动 revival}

同一个势 \eqnrefs{eq:linear-molecule-alignment-potential} 在不同脉冲时间尺度下产生两种很不相同的物理。对线性转子，场自由能谱为

\[
E_J=B_eJ(J+1),
\]

因此典型转动时间尺度为 \(\hbar/B_e\)，完整量子 revival 时间为

\begin{equation}
T_{\rm rev}=\frac{\pi\hbar}{B_e}.
\label{eq:rotational-revival-time}
\end{equation}

若脉冲包络变化时间 \(\tau_p\) 满足

\[
\tau_p\gg\frac{\hbar}{|E_n-E_m|}
\]

并且更严格的绝热条件

\begin{equation}
\epsilon_{nm}^{\rm ad}
\equiv
\frac{\hbar|\langle n(t)|\dot H(t)|m(t)\rangle|}
{|E_n(t)-E_m(t)|^2}
\ll1
\label{eq:rotational-adiabatic-parameter}
\end{equation}

对所有重要耦合态成立，分子会跟随瞬时 pendular eigenstate，称 adiabatic alignment。若场也绝热关闭，系统大致回到原来的无场转动态，强 alignment 不会长期保留。

相反，当

\[
\tau_p\ll\frac{\hbar}{B_e},
\]

脉冲期间自由转动来不及显著演化，可以把激光看成一个瞬时“kick”。忽略脉冲内的 \(H_{\rm rot}\) 后，演化算符近似为

\begin{equation}
U_{\rm kick}
\simeq
\exp\!\left[
\ii P\cos^2\theta
\right],
\qquad
P=\frac{\Delta\alpha}{4\hbar}
\int_{-\infty}^{\infty}E_0^2(t)\,\dd t.
\label{eq:alignment-kick-strength}
\end{equation}

无量纲 \(P\) 是 impulsive alignment 的 kick strength。脉冲结束后，多个 \(J\) 态按照 \(E_J\) 自由积累相位；正是这些由 \eqnrefs{eq:alignment-selection-rule} 逐级建立的 \(\Delta J=\pm2\) 相干在 \(T_{\rm rev}\) 的整数或分数倍附近重新相位匹配，产生 field-free alignment revival。


\begin{derivation}[从短脉冲传播子到完整转动 revival]
脉冲期间的精确传播子为
\[
U(t_+,t_-)
=\mathcal T\exp\left[-\frac{\ii}{\hbar}
\int_{t_-}^{t_+}\bigl(H_{\rm rot}+V(t)\bigr)\,\dd t\right].
\]
若脉冲只显著占据一组有限的转动态，并且该子空间内满足
\[
\epsilon_{\rm imp}
\equiv
\frac{\tau_p\,\Delta E_{\rm rot}^{(\rm acc)}}{\hbar}\ll1,
\]
则在脉冲宽度内自由转动造成的相对相位可以作为受控小量。Magnus 展开的首项给出
\[
U(t_+,t_-)
=\exp\left[-\frac{\ii}{\hbar}\int V(t)\,\dd t\right]
+O(\epsilon_{\rm imp}),
\]
而 \(J\)-无关的标量 AC Stark 项只贡献整体相位。代入\eqnrefs{eq:linear-molecule-alignment-potential}便得到
\[
U_{\rm kick}
=\exp\!\left(\ii P\cos^2\theta\right)
+O(\epsilon_{\rm imp}).
\]

脉冲后线性刚性转子自由演化。对任意整数 \(J\)，
\[
U_0(T_{\rm rev})|JM\rangle
=\exp\!\left[-\frac{\ii B_eJ(J+1)}{\hbar}
\frac{\pi\hbar}{B_e}\right]|JM\rangle.
\]
由于相邻整数的乘积 \(J(J+1)\) 必为偶数，
\[
\exp[-\ii\pi J(J+1)]=1,
\]
所以所有被 kick 激发的 \(J\) 分量在
\[
T_{\rm rev}=\frac{\pi\hbar}{B_e}
\]
处同时恢复相位，得到完整 revival。分数 revival 则来自同一二次谱相位在 \(T_{\rm rev}/q\) 处按同余类重新组织，而不是额外引入的新动力学机制。

\textbf{物理意义。}分子转动 revival 是量子波包在离散转动能级 $E_J=BJ(J+1)$ 上演化的干涉效应。经典上刚性转子波包会迅速弥散（不同 $J$ 相速度不同），但在 $T_{\rm rev}=\pi\hbar/B$ 后所有相位重整，波包复原。实例：$\mathrm{N}_2$ 的 $B=1.998\,\si{\cm^{-1}}$，$T_{\rm rev}\approx8.4\,\si{\ps}$，在飞秒激光脉冲激发后用 pump-probe 实验直接观测到周期性波包复原。分数 revival 在 $T_{\rm rev}/4$ 处出现四簇波包，是 $E_J$ 二次依赖的直接量子签名。
\end{derivation}

\begin{note}[Alignment versus orientation]
\eqnrefs{eq:linear-molecule-alignment-potential} 只含 \(\cos^2\theta\)，在 \(\hat{\mathbf n}\to-\hat{\mathbf n}\) 下不变，所以它只能排列分子轴而不能区分"头"和"尾"。真正的 orientation 需要奇宇称耦合，例如永久偶极与静电场的 \(-\boldsymbol\mu\cdot\mathbf E\) 或合适双色场产生的奇阶有效相互作用。
\end{note}

\section{转动光谱实验与分子结构测定}\label{sec:rotation-experiments}

本节给出转动光谱的典型实验参数、不对称陀螺的数值实例，以及激光 alignment 的实验参数。

接下来考察双原子分子转动常数与纯转动光谱。

纯转动光谱（微波区）由 \(\Delta J=1\) 电偶极跃迁给出，频率
\[
\tilde\nu(J\to J+1)=2\tilde B_e(J+1)-4\tilde D_e(J+1)^3
\]
永久偶极矩 \(\mu\ne0\) 是纯转动光谱活性的必要条件：同核双原子（H\(_2\)、N\(_2\)、O\(_2\)）无纯转动红外/微波谱。

\begin{center}
\begin{tabular}{lcccc}
\hline
分子 & \(\tilde B_e\) (cm\(^{-1}\)) & \(\mu\) (D) & \(T_{\rm rev}=\pi/(2\tilde B_ec)\) (ps) & 跃迁频率 \(J=0\to1\) (GHz) \\
\hline
HCl & 10.593 & 1.08 & 0.79 & 634.8 \\
CO & 1.931 & 0.112 & 4.33 & 115.8 \\
NO & 1.702 & 0.153 & 4.91 & 102.2 \\
OCS & 0.203 & 0.715 & 41.2 & 12.2 \\
HCN & 1.478 & 3.00 & 5.65 & 88.7 \\
\hline
\end{tabular}
\end{center}

\begin{exbox}[CO 的纯转动光谱]
CO 的 \(J=0\to1\) 跃迁在 115.8\,GHz（\(\lambda\approx2.59\,\mathrm{mm}\)），是天体化学中探测星际 CO 分子的标志性谱线。前几条线：\(J=1\to2\) 在 231.6\,GHz，\(J=2\to3\) 在 347.4\,GHz——等间距 \(2\tilde B_ec\approx115.8\,\mathrm{GHz}\) 是刚性转子的特征。离心畸变使间距随 \(J\) 略减：\(J=10\to11\) 的间距为 115.5\,GHz，偏差 \(\sim0.3\%\)。
\end{exbox}

接下来考察不对称陀螺：水分子与甲醛。

水分子的主轴转动常数 \(A_{\rm rot}/(hc)=27.877\,\mathrm{cm}^{-1}\)，\(B_{\rm rot}/(hc)=14.512\,\mathrm{cm}^{-1}\)，\(C_{\rm rot}/(hc)=9.285\,\mathrm{cm}^{-1}\)，Ray 参数
\[
\kappa_{\rm R}=\frac{2B_{\rm rot}-A_{\rm rot}-C_{\rm rot}}{A_{\rm rot}-C_{\rm rot}}
=\frac{29.024-27.877-9.285}{27.877-9.285}\approx-0.437
\]
中等不对称，\(K\) 混合显著。

\begin{exbox}[\(J=1\) 水分子能级]
不对称陀螺 \(J=1\) 有 3 个能级（\(M\) 简并不计）。在 \(|JKM\rangle\) 基中，\(K=0,\pm1\) 分为 \(K\) 偶（\(K=0\)）和 \(K\) 奇（\(K=\pm1\)）两个块。\(K\) 偶块：\(E_{1,0}=\frac12(A_{\rm rot}+B_{\rm rot})+C_{\rm rot}-\frac12(A_{\rm rot}+B_{\rm rot})=C_{\rm rot}=9.285\,\mathrm{cm}^{-1}\)。\(K\) 奇块是 \(2\times2\) 矩阵，本征值 \(\frac12(A_{\rm rot}+B_{\rm rot})\pm\frac12\sqrt{(A_{\rm rot}-B_{\rm rot})^2+4\cdot0}\)：\(\frac12(A_{\rm rot}+B_{\rm rot})=21.195\) 和 \(A_{\rm rot}\) 或 \(B_{\rm rot}\) 的线性组合。精确对角化给出 23.398 和 14.512\,cm\(^{-1}\)。三个能级 9.285、14.512、23.398\,cm\(^{-1}\) 与实验值 9.285、14.512、23.398\,cm\(^{-1}\) 完全吻合。
\end{exbox}

接下来考察激光 alignment 实验参数。

\eqnrefs{eq:alignment-kick-strength} 的 kick strength \(P\) 取决于极化率各向异性 \(\Delta\alpha\) 和脉冲能量。典型参数：

\begin{center}
\begin{tabular}{lcccc}
\hline
分子 & \(\Delta\alpha\) (\AA\(^3\)) & \(B_e/h\) (GHz) & \(T_{\rm rev}\) (ps) & 典型 \(P\) \\
\hline
N\(_2\) & 0.69 & 59.9 & 8.4 & 5--10 \\
O\(_2\) & 0.60 & 43.4 & 11.5 & 4--8 \\
CO\(_2\) & 1.85 & 11.7 & 42.8 & 10--20 \\
I\(_2\) & 6.0 & 1.12 & 448 & 50--100 \\
\hline
\end{tabular}
\end{center}

\begin{exbox}[\(I_2\) 的转动 revival]
I\(_2\) 的 \(T_{\rm rev}\approx448\,\mathrm{ps}\) 是实验上最容易观测的 revival 之一。Stapelfeldt 实验用 800\,nm、100\,fs 脉冲（\(\tau_p\ll\hbar/B_e\approx140\,\mathrm{ps}\)）激发，kick strength \(P\approx60\)。在 \(T_{\rm rev}/4\approx112\,\mathrm{ps}\) 观测到 4 簇波包（quarter revival），在 \(T_{\rm rev}/2\approx224\,\mathrm{ps}\) 观测到 2 簇（half revival），在 \(T_{\rm rev}\approx448\,\mathrm{ps}\) 完全复原。分数 revival 的存在直接验证了 \(E_J\propto J(J+1)\) 的二次依赖——若能级是线性的，revival 不会出现分数结构。
\end{exbox}

\begin{derivation}[不对称陀螺的 Wang 变换对角化与能级对称分类，即\eqnrefs{eq:wang-block-diagonalization}]
\emph{策略路线图。}首先写出不对称陀螺 Hamiltonian 在对称陀螺基 \(|JKM\rangle\) 中的矩阵表示并识别 \(K\to K\pm2\) 混合结构；其次构造 Wang 变换将 \(|JK\rangle\) 与 \(|J,-K\rangle\) 组合成确定宇称态；再次证明 Wang 基下 Hamiltonian 分解为四个不耦合子块；然后对 \(J=1,2\) 显式对角化并标注能级对称类；最后用 Ray 不对称参数 \(\kappa\) 连接长陀螺与扁陀螺极限。

\emph{第一步：不对称陀螺矩阵元。}不对称陀螺 Hamiltonian
\begin{equation}
H=A_{\rm rot}J_a^2+B_{\rm rot}J_b^2+C_{\rm rot}J_c^2,
\qquad
A_{\rm rot}\ge B_{\rm rot}\ge C_{\rm rot},
\label{eq:asymmetric-top-hamiltonian-wang}
\end{equation}
即\eqnrefs{eq:asymmetric-top-hamiltonian}。改写为对称陀螺部分加微扰：
\begin{equation*}
H=\bar B J^2+(A_{\rm rot}-C_{\rm rot})\!\left[\frac{1+\kappa}{2}J_a^2+\frac{1-\kappa}{2}J_c^2\right],
\qquad
\bar B=\frac{A_{\rm rot}+C_{\rm rot}}{2},
\qquad
\kappa=\frac{2B_{\rm rot}-A_{\rm rot}-C_{\rm rot}}{A_{\rm rot}-C_{\rm rot}},
\end{equation*}
\(\kappa\in[-1,+1]\) 为 Ray 不对称参数（\(\kappa=-1\) 长陀螺，\(\kappa=+1\) 扁陀螺）。在 \(|JKM\rangle\) 基中，\(J_a^2\) 和 \(J_c^2\) 的矩阵元由\eqnrefs{eq:asymmetric-top-diagonal-element} 和\eqnrefs{eq:asymmetric-top-offdiagonal-element} 给出：对角项来自 \(J_c^2\)，非对角项来自 \(J_a^2\) 的 \(K\to K\pm2\) 跃迁
\begin{equation}
\langle J,K\pm2|H|J,K\rangle
=\frac{A_{\rm rot}-C_{\rm rot}}{4}
\bigl[J(J+1)-K(K\pm1)\bigr]^{1/2}
\bigl[J(J+1)-(K\pm1)(K\pm2)\bigr]^{1/2}.
\label{eq:asymmetric-top-k-mixing-wang}
\end{equation}

\emph{第二步：Wang 变换。}其次注意\eqnrefs{eq:asymmetric-top-k-mixing-wang} 只耦合 \(K\) 与 \(K\pm2\)，故 \(K\) 的奇偶性守恒。定义 Wang 组合态
\begin{equation}
|J,K,\sigma\rangle=\frac{1}{\sqrt2}\bigl(|J,K\rangle+\sigma\,|J,-K\rangle\bigr),
\qquad
K>0,\ \sigma=\pm1,
\label{eq:wang-transform}
\end{equation}
以及 \(|J,0\rangle\)（\(K=0\) 自成一组，\(\sigma=+1\)）。Wang 变换是 \(K\) 空间中的离散 Fourier 变换，将 body-fixed \(C_2\) 绕 \(b\) 轴旋转的对称性对角化。

\emph{第三步：四块分解。}再次在 Wang 基下计算矩阵元。\(K\to K+2\) 的耦合在 \(|J,K,\sigma\rangle\) 与 \(|J,K+2,\sigma'\rangle\) 间只在 \(\sigma=\sigma'\) 时非零——Wang 宇称 \(\sigma\) 守恒。同时 \(K\) 奇偶性也守恒。故 Hamiltonian 分解为四个不耦合子块
\begin{equation}
H=H_{E^+}\oplus H_{E^-}\oplus H_{O^+}\oplus H_{O^-},
\label{eq:wang-block-diagonalization}
\end{equation}
即\eqnrefs{eq:wang-block-diagonalization}，其中 \(E/O\) 标记 \(K\) 偶/奇，\(\pm\) 标记 Wang 宇称 \(\sigma=\pm1\)。每个子块维数约 \(J/4\)，对角化效率提升 16 倍。

\emph{第四步：\(J=1\) 显式对角化。}然后对 \(J=1\)，\(K=0,\pm1\)。Wang 基：\(|1,0,+\rangle\)（\(E^+\) 块，1 维），\(|1,1,+\rangle\)（\(O^+\) 块，1 维），\(|1,1,-\rangle\)（\(O^-\) 块，1 维）。\(E^+\) 块：\(K=0\) 无 \(K\pm2\) 耦合，
\begin{equation*}
E_{E^+}=\bar B\cdot2+(A_{\rm rot}-C_{\rm rot})\frac{1-\kappa}{2}\cdot0=C_{\rm rot}.
\end{equation*}
\(O^+\) 块：\(K=1\) 对角项 \(\bar B\cdot2+(A_{\rm rot}-C_{\rm rot})\frac{1-\kappa}{2}\)，化简为 \(B_{\rm rot}\)。\(O^-\) 块：对角项 \(\bar B\cdot2+(A_{\rm rot}-C_{\rm rot})\frac{1+\kappa}{2}\)，化简为 \(A_{\rm rot}\)。三个能级
\begin{equation*}
E_{J=1}\in\{A_{\rm rot},\ B_{\rm rot},\ C_{\rm rot}\},
\end{equation*}
分别属于 \(O^-\)、\(O^+\)、\(E^+\) 对称类。长陀螺极限 \(\kappa\to-1\)（\(B_{\rm rot}\to C_{\rm rot}\)）：\(E_{O^+}=E_{E^+}=C_{\rm rot}\) 简并，\(E_{O^-}=A_{\rm rot}\) 分离——恢复长陀螺 \(K=0\) 与 \(K=1\) 能级结构。

\emph{第五步：\(J=2\) 与简并解除。}最后对 \(J=2\)，\(K=0,\pm1,\pm2\)。\(E^+\) 块含 \(|2,0,+\rangle\)、\(|2,2,+\rangle\)（2 维），\(O^-\) 块含 \(|2,1,-\rangle\)（1 维），\(O^+\) 块含 \(|2,1,+\rangle\)（1 维），\(E^-\) 块含 \(|2,2,-\rangle\)（1 维）。\(E^+\) 块的 \(2\times2\) 矩阵
\begin{equation*}
H_{E^+}^{(J=2)}
=\begin{pmatrix}
6\bar B & \sqrt{3}(A_{\rm rot}-C_{\rm rot})\\
\sqrt{3}(A_{\rm rot}-C_{\rm rot}) & 6\bar B+4(A_{\rm rot}-C_{\rm rot})\frac{1-\kappa}{2}
\end{pmatrix},
\end{equation*}
本征值 \(E_{E^+}^{\pm}=6\bar B+(A_{\rm rot}-C_{\rm rot})(1-\kappa)\pm\sqrt{3(A_{\rm rot}-C_{\rm rot})^2+(A_{\rm rot}-C_{\rm rot})^2(1-\kappa)^2/4}\)，在 \(\kappa=0\)（最大不对称）给出最大能级分裂。水分子 \(\kappa\approx-0.437\) 处于中等不对称，\(J=2\) 的 5 个能级（2 个 \(E^+\) + 1 个 \(E^-\) + 1 个 \(O^+\) + 1 个 \(O^-\)）全部分裂，与\eqnrefs{eq:asymmetric-top-k-mixing} 的微扰分析一致。


\emph{[三维转子的定向排列与超极化矩]} 激光诱导分子排列利用超短脉冲的各向异性极化将随机取向的分子定向排列，其理论基于三维转子的量子动力学。分四步建立。

\emph{第一步：刚性转子 Hamiltonian。}刚性线性转子 \(H_r=B\boldsymbol J^2\)（\(B=\hbar^2/(2I)\) 转动常数），本征态 \(|J,M\rangle\)（\(J=0,1,\ldots\)，\(M=-J,\ldots,J\)），能级 \(E_J=BJ(J+1)\)。与外场耦合 \(H_{\rm int}=-\boldsymbol\mu\cdot\boldsymbol E-\frac{1}{2}\boldsymbol\alpha\cdot\boldsymbol E\boldsymbol E\)（偶极 + 极化）。对非极性分子（\(\mu=0\)），极化各向异性 \(\Delta\alpha=\alpha_\parallel-\alpha_\perp\) 给出 \(H_{\rm int}=-\frac{\Delta\alpha E^2}{2}\cos^2\theta\)（\(\theta\) 是分子轴与场方向夹角）。

\emph{第二步：脉冲排列与绝热条件。}超短脉冲（\(\tau\ll1/B\)，非绝热）将转子从 \(J=0\) 态"踢"到 \(J\) 态叠加：\(|\psi\rangle=\sum_Jc_J|J,0\rangle\)。脉冲后的自由演化 \(|\psi(t)\rangle=\sum_Jc_J\ee^{-iBJ(J+1)t}|J,0\rangle\) 在 revival 时间 \(t_{\rm rev}=\pi/(2B)\) 回到初始态，在 \(t_{\rm rev}/2\) 达到最大排列（\(\langle\cos^2\theta\rangle\) 最大）。绝热条件 \(\tau\gg1/B\) 下转子跟随瞬时本征态，排列度由势阱深度 \(\Delta\alpha E^2/(2B)\) 决定。

\emph{第三步：排列度与定向度。}排列度 \(\langle\cos^2\theta\rangle\)（\(\frac{1}{3}\) 各向同性，\(1\) 完全排列）和定向度 \(\langle\cos\theta\rangle\)（\(0\) 无定向，\(1\) 完全定向）是两个独立量度。线性极化光只排列（\(\langle\cos\theta\rangle=0\)，\(\theta\) 和 \(\pi-\theta\) 等权），椭圆偏振光可定向（\(\langle\cos\theta\rangle\neq0\)）。完全定向需要打破 \(\theta\to\pi-\theta\) 对称，如双色光（\(\omega+2\omega\)）或静电场组合。

\emph{第四步：三维排列与复杂分子。}非线性分子（三个转动轴 \(A,B,C\)）的排列涉及三维 Euler 角 \((\phi,\theta,\chi)\)，转动 Hamiltonian \(H_r=AJ_a^2+BJ_b^2+CJ_c^2\)（不对称转子）或 \(H_r=B\boldsymbol J^2+(A-B)J_a^2\)（对称转子）。极化耦合 \(H_{\rm int}=-\frac{1}{2}\sum_{ij}\alpha_{ij}E_iE_j\) 含极化张量 \(\alpha_{ij}\) 的全部分量。三维排列的目标是同时固定三个 Euler 角——需要三维各向异性场（如三正交偏振光）。三维排列度用 \(\langle R_{ij}\rangle\)（方向余弦矩阵元）量度，在分子碰撞和化学反应控制中应用——排列的分子有方向性反应截面，是立体动力学的基础。
\end{derivation}
